\section{System: \Weld}
\label{sec:system}

\subsection{Architecture Overview}
\label{sec:system:overview}

\Weld is an open-source, full-stack developer workspace targeting teams of
two to ten engineers. The web frontend is built with SvelteKit 2 (TypeScript,
Tailwind CSS 4); the backend is an Axum 0.8 server (Rust) backed by PostgreSQL
16 and Redis. A WebSocket hub broadcasts real-time board and chat updates to all
connected clients of a team. All study-mode paths share the same production
application binary, activated by an \texttt{appConfig.studyMode} flag.

\Weld exposes three collaboration surfaces that correspond exactly to the three
context sources manipulated across the agent conditions:

\begin{description}
  \item[Kanban Board.]
    Columns and cards with fractional-indexed positions, real-time drag-and-drop
    reordering, label and priority assignment, and GitHub issue/PR links. Cards
    carry a \texttt{github\_issue} and \texttt{github\_pr} field that are
    populated either by the user or by the agent's \texttt{link\_card\_to\_github}
    tool. Cross-references to chat messages render inline as
    \texttt{[[card:uuid]]} tokens that resolve to card titles on hover.

  \item[Team Chat.]
    Channel-based messaging with threading, Markdown rendering, @mention
    resolution, and the same \texttt{[[card:uuid]]} cross-reference syntax used
    on the board. Messages are delivered over WebSocket and persisted
    server-side; the agent's \texttt{post\_message} tool drafts a message that
    the participant must review and confirm before it is sent.

  \item[GitHub Integration.]
    A GitHub App installation synchronizes open issues and pull requests into
    \Weld's read-only GitHub view. Two-way linking (card $\leftrightarrow$
    issue, card $\leftrightarrow$ PR) is supported both manually and via the
    agent's \texttt{link\_card\_to\_github} tool.
\end{description}

\subsection{Study Mode and Workspace Isolation}
\label{sec:system:study}

Each study participant is provisioned in a single atomic database transaction
via the \texttt{provision\_participant} function. The transaction creates:
a dedicated user account, a team with four pre-seeded synthetic teammates
(Sarah Chen, Alex Rivera, Priya Nair, and Marcus Webb), a PetTracker
project, three board columns (Backlog, In Progress, Done), 15--20 cards with
realistic descriptions and priorities, two channels (\#general and
\#pettracker) with pre-seeded message history, and GitHub fixture data.
All timestamps in the seeded content are stored as negative offsets from the
provisioning moment, so the workspace always appears to have been active for
the same relative period regardless of when the participant joins.

\paragraph{Access codes and resumption.}
Participants authenticate via single-use access codes rather than email/password
registration. If a participant closes the browser mid-session, they can return
to the same code on any device and resume from their saved stage and task index
(\texttt{resume\_stage}, \texttt{resume\_task\_index}). Progress is persisted by
\texttt{update\_progress(pool, user\_id, stage, task\_index)}.

\paragraph{Condition assignment.}
The \texttt{pick\_agent\_condition} function implements balanced permuted-block
randomization with block size 3 and a random tie-break on equal counts, ensuring
that no condition falls behind the others as enrollment proceeds while avoiding
deterministic C1$\to$C2$\to$C3 cycling that could confound condition with
enrollment order.

\paragraph{Database isolation.}
Every query in study mode is scoped by \texttt{team\_id} and
\texttt{project\_id} in the SQL \texttt{WHERE} clause. No cross-participant
query paths exist; the study database schema enforces this at the row level.

\subsection{The Context Gate}
\label{sec:system:gate}

The \ContextGate is a Rust module (\texttt{services/study/agent/gate.rs}) that
acts as the sole chokepoint for assembling the agent's context before each LLM
call. Its design centers on two structures: \texttt{ContextRules}, which encodes
the per-condition allowlist, and \texttt{ScopedContext}, the assembled result.

\paragraph{Rules and allowlists.}
\texttt{rules\_for(condition)} returns a \texttt{ContextRules} value specifying
which of the three sources (chat, board, GitHub) may be queried. The rules for
each condition are:

\begin{center}
\begin{tabular}{lccc}
  \toprule
  Condition & Chat & Board & GitHub \\
  \midrule
  \CI\ (advisory)     & \checkmark & -- & -- \\
  \CII\ (action)      & \checkmark & \checkmark & -- \\
  \CIII\ (full)       & \checkmark & \checkmark & \checkmark \\
  \bottomrule
\end{tabular}
\end{center}

The gate's \texttt{assemble(pool, ctx)} function queries \emph{only} the allowed
sources; the database calls for disallowed sources are never made. The resulting
\texttt{ScopedContext} carries a \texttt{sources\_included} field listing exactly
which sources were consulted, a \texttt{chat: Vec<ChatLine>} field, an
\texttt{board: Option<Vec<CardLine>>}, and a \texttt{github: Option<Vec<GithubItem>>}.
Disallowed fields are \texttt{None}.

\paragraph{Proof-of-scoping snapshot.}
Before invoking the LLM, \texttt{run\_turn} calls
\texttt{ScopedContext::snapshot()}, which serialises the assembled context (minus
any fields that should not appear in logs) and records it as an
\texttt{agent\_context\_sent} telemetry event. This creates an immutable per-turn
audit record proving that the context assembly respected the assigned condition.

\paragraph{Double-gated tool executor.}
The gate also governs which tools are available to the model.
\texttt{allowed\_tool\_categories(condition)} returns a static slice of
\texttt{ToolCategory} variants: \CI\ receives \texttt{[]}, \CII\ receives
\texttt{[BoardWrite, ChatPost]}, and \CIII\ receives all three categories
(\texttt{[BoardWrite, ChatPost, GithubLink]}). This list determines which tool
definitions are sent to the model in the OpenAI request (first gate). The tool
executor (\texttt{tools::execute}) independently re-validates the tool name
against the condition's allowed categories before executing any action; any
call that slips past the first gate is rejected and logged as
\texttt{agent\_tool\_rejected} (second gate). Table~\ref{tab:tools} lists the
five tools and their condition availability.

\begin{table}[h]
  \caption{Agent tools by condition availability. \checkmark~= available,
    --~= not offered to model and rejected at executor.}
  \label{tab:tools}
  \begin{tabular}{llccc}
    \toprule
    Tool name & Action & \CI & \CII & \CIII \\
    \midrule
    \texttt{create\_card}        & Create a board card           & -- & \checkmark & \checkmark \\
    \texttt{move\_card}          & Move a card to another column & -- & \checkmark & \checkmark \\
    \texttt{update\_card}        & Edit card priority / assignee / description & -- & \checkmark & \checkmark \\
    \texttt{post\_message}       & Draft a chat message for review & -- & \checkmark & \checkmark \\
    \texttt{link\_card\_to\_github} & Link a card to an issue or PR & -- & -- & \checkmark \\
    \bottomrule
  \end{tabular}
\end{table}

\paragraph{Turn budget.}
Each agent invocation is bounded by two hard limits: \texttt{MAX\_TURNS~=~5}
model$\leftrightarrow$tool round-trips per invocation, and
\texttt{MAX\_AGENT\_MESSAGES~=~50} total agent messages per participant across
the entire session. These limits prevent runaway loops and bound infrastructure
cost.

\begin{figure}[h]
  \centering
  % TODO: Replace with final architecture diagram
  \fbox{\rule{0pt}{5cm}\rule{0.85\linewidth}{0pt}}
  \caption{The \ContextGate architecture. The \texttt{assemble()} call queries only
    allowed sources for the participant's assigned condition and emits a
    \texttt{agent\_context\_sent} snapshot. The double-gated tool executor
    re-validates tool permissions independently at execution time, logging
    violations as \texttt{agent\_tool\_rejected} events.}
  \label{fig:gate}
\end{figure}

\subsection{Agent Turn Orchestration}
\label{sec:system:agent}

The \texttt{run\_turn} function orchestrates each agent invocation. After
assembling the \texttt{ScopedContext} and recording the
\texttt{agent\_context\_sent} event, it loads prior turns from the
\texttt{study\_agent\_turns} table to maintain conversation continuity across
the session. The function then enters a loop bounded by \texttt{MAX\_TURNS},
alternating between LLM calls (with condition-appropriate tool definitions from
\texttt{openai\_tool\_defs(condition)}) and tool execution calls.

\paragraph{Condition-specific system prompt.}
The system prompt varies by condition. For \CI, it instructs the model to
``provide guidance but not suggest or execute actions,'' constraining it to an
advisory role even though no tools are offered. For \CII\ and \CIII, the prompt
frames the agent as ``a co-worker in the team workspace'' that can take actions
on behalf of the participant. This framing difference is intentional: it matches
participant expectations to the agent's actual capability.

\paragraph{Chat message drafting.}
The \texttt{post\_message} tool does not immediately post a message; instead,
the tool executor returns a \texttt{ChatDraft} to the client, which renders it
as a review panel. The participant must explicitly confirm before the message is
sent. This keeps human oversight over public communication while still reducing
the mechanical effort of composing and locating the correct channel.

\subsection{Telemetry Infrastructure}
\label{sec:system:telemetry}

\Weld records 43 discrete event types (Table~\ref{tab:events}). Client events
are batched and submitted to a single endpoint that validates each event type
against a server-side allowlist (\texttt{ALLOWED\_EVENT\_TYPES}) before
persistence; events not on the allowlist are silently dropped. This
\emph{allowlist-by-construction} approach means that even a modified client
cannot introduce new event types, preventing accidental capture of keystrokes,
scroll positions, or other data outside the study protocol.

Server-side events (agent lifecycle, context snapshots, tool outcomes) are
recorded via \texttt{record\_server\_event} directly within the Axum handlers,
bypassing the client batch path and therefore unaffected by client connectivity
or manipulation.

\begin{table}[h]
  \caption{Telemetry event categories and types recorded by \Weld.}
  \label{tab:events}
  \small
  \begin{tabular}{lp{8.5cm}}
    \toprule
    Category & Event types \\
    \midrule
    Session \& phase
      & \texttt{session\_start}, \texttt{session\_end},
        \texttt{session\_abandoned}, \texttt{session\_resumed},
        \texttt{phase\_started}, \texttt{phase\_ended},
        \texttt{consent\_given} \\
    Survey
      & \texttt{survey\_started}, \texttt{survey\_completed} \\
    Task lifecycle
      & \texttt{task\_shown}, \texttt{task\_completed}, \texttt{task\_skipped},
        \texttt{task\_final\_state}, \texttt{task\_rating},
        \texttt{instruction\_reopened} \\
    Board
      & \texttt{card\_created}, \texttt{card\_opened}, \texttt{card\_edited},
        \texttt{card\_moved}, \texttt{card\_assigned}, \texttt{card\_labeled},
        \texttt{card\_priority\_set}, \texttt{card\_deleted},
        \texttt{card\_searched}, \texttt{board\_filtered} \\
    Chat
      & \texttt{message\_sent}, \texttt{message\_edited},
        \texttt{message\_deleted}, \texttt{channel\_switched},
        \texttt{thread\_opened}, \texttt{mention\_clicked},
        \texttt{link\_clicked} \\
    Cross-surface
      & \texttt{card\_referenced\_from\_chat},
        \texttt{chat\_referenced\_from\_card},
        \texttt{pr\_referenced\_from\_card},
        \texttt{card\_referenced\_from\_pr},
        \texttt{issue\_referenced\_from\_card} \\
    GitHub
      & \texttt{github\_view\_opened}, \texttt{pr\_opened},
        \texttt{issue\_opened}, \texttt{pr\_linked\_to\_card},
        \texttt{issue\_linked\_to\_card} \\
    Navigation
      & \texttt{view\_changed}, \texttt{search\_used} \\
    Agent
      & \texttt{agent\_invoked}, \texttt{agent\_context\_sent},
        \texttt{agent\_response}, \texttt{agent\_tool\_call},
        \texttt{agent\_tool\_rejected}, \texttt{agent\_suggestion\_accepted},
        \texttt{agent\_suggestion\_rejected}, \texttt{agent\_action\_undone},
        \texttt{agent\_output\_edited}, \texttt{agent\_error},
        \texttt{agent\_panel\_opened}, \texttt{agent\_panel\_collapsed} \\
    \bottomrule
  \end{tabular}
\end{table}
